\documentclass[11pt]{article}
\usepackage[margin=1in]{geometry}
\usepackage{amsmath,amssymb,amsthm,graphicx,booktabs}
\usepackage[hidelinks]{hyperref}
\newtheorem{theorem}{Theorem}
\newtheorem{lemma}{Lemma}
\newtheorem{corollary}{Corollary}
\newcommand{\R}{\mathbb R}
\newcommand{\Z}{\mathbb Z}
\newcommand{\D}{D_+}
\newcommand{\tht}{\theta}
\setlength{\parindent}{0pt}
\setlength{\parskip}{0.62em}

\title{A Counterexample to the Proposed Stability Conditions for Gaussian Lattice Sums}
\author{Run \texttt{fa\_banach\_001}, lane 18}
\date{11 August 2026}

\begin{document}
\maketitle

\textbf{Status.} Candidate full counterexample to the Appendix~C conjectured
sufficient conditions in arXiv:2110.06008, likely valid, for human review.

\section{Source conjecture}

For $\tau=x+iy$ in
\[
 \D=\{(x,y):0\leq x\leq\tfrac12,\ x^2+y^2\geq1,\ y>0\},
\]
B\'etermin, Faulhuber, and Steinerberger parametrize a unit-area lattice and
define
\begin{equation}\label{eq:theta}
 \tht_{x,y}(z;\alpha)=
 \sum_{k,l\in\Z}\exp\!\left[-\frac{\pi\alpha}{y}
 \left((k+z_1)^2+2x(k+z_1)(l+z_2)
 +(x^2+y^2)(l+z_2)^2\right)\right].
\end{equation}
The hexagonal lattice has $(x,y)=(1/2,\sqrt3/2)$ and deep hole
$z^-=(1/3,1/3)$.  Appendix~C proposes that
\begin{equation}\label{eq:conj}
 \tht_{1/2,\sqrt3/2}(z^-;\alpha)
 \geq \tht_{x,y}(\widetilde z(x,y);\alpha)
 \qquad(\alpha>0)
\end{equation}
at least when the selector obeys
\begin{equation}\label{eq:conditions}
 \widetilde z_1+x\widetilde z_2=\tfrac12,\qquad
 \partial_x\widetilde z_2\leq0,\qquad
 \widetilde z(\tfrac12,y)=a(\tfrac12,y),\qquad
 \widetilde z(0,y)\in[\tfrac14,\tfrac12],
\end{equation}
and tends to $z^-$ with the lattice \cite[Appendix~C]{BFS}.

\begin{center}
\includegraphics[width=0.94\textwidth]{figures/open_problem_crop.png}
\end{center}

The last membership in \eqref{eq:conditions} is dimensionally ambiguous for a
vector-valued selector; the construction below has
$\widetilde z(0,y)=(1/2,1/2)$, so it satisfies both the natural reading as a
condition on $\widetilde z_2$ and the coordinatewise reading.

\section{Claimed counterexample}

\begin{theorem}\label{thm:main}
Define on all of $\D$
\begin{equation}\label{eq:selector}
 \widetilde z_2(x,y)=\frac12-\frac{[4x(1-x)]^{100}}{8y^2},
 \qquad
 \widetilde z_1(x,y)=\frac12-x\widetilde z_2(x,y).
\end{equation}
This selector satisfies every condition in \eqref{eq:conditions}, tends to
$(1/3,1/3)$ at the hexagonal lattice, and also satisfies the stronger boundary
condition
\[
 \partial_x\widetilde z_2(\tfrac12,y)=0.
\]
Nevertheless, at
\begin{equation}\label{eq:test}
 x=\frac9{20},\qquad y=\frac{\sqrt{319}}{20},\qquad \alpha=1,
\end{equation}
one has the strict reverse inequality
\[
 \tht_{x,y}(\widetilde z(x,y);1)
 >\tht_{1/2,\sqrt3/2}((1/3,1/3);1).
\]
Hence the proposed sufficient conditions do not imply \eqref{eq:conj}, even
after adding the boundary-criticality property discussed immediately after
the conjecture.
\end{theorem}

\section{Idea of the proof}

The endpoint constraints control how far the selector can move but do not
control how rapidly it moves as $x$ leaves the hexagonal boundary.  The factor
$[4x(1-x)]^{100}$ equals one at $x=1/2$, has zero derivative there, yet drops
quickly just inside the domain.  Thus \eqref{eq:selector} remains perfectly
smooth and monotone while moving much closer to $(1/2,1/2)$ than the canonical
selectors considered in the source paper.  That displacement raises the theta
value enough to cross the hexagonal benchmark.

To make the crossing rigorous, no infinite numerical sum is trusted.  We keep
only a $5\times5$ square of positive terms for the candidate, giving a lower
bound.  For the hexagonal sum, its quadratic form dominates a separable one;
the omitted complement is then bounded by two elementary one-dimensional
Gaussian geometric tails.  Directed interval arithmetic leaves a gap above
$0.003847$.

\section{Proof}

Put $G(x)=4x(1-x)$.  On $0\leq x\leq1/2$ one has $0\leq G\leq1$ and
$G'(x)=4(1-2x)\geq0$.  The line condition in \eqref{eq:conditions} is built
into \eqref{eq:selector}, while
\[
 \partial_x\widetilde z_2
 =-\frac{100}{8y^2}G(x)^{99}G'(x)\leq0.
\]
At $x=0$, $G=0$ and $\widetilde z=(1/2,1/2)$.  At $x=1/2$, $G=1$ and
\[
 \widetilde z_2=\frac12-\frac1{8y^2},\qquad
 \widetilde z_1=\frac14+\frac1{16y^2},
\]
which are exactly the source paper's circumcenter coordinates
$a(1/2,y)$.  Moreover $G'(1/2)=0$, proving boundary criticality.  Continuity
at $(1/2,\sqrt3/2)$ gives $\widetilde z\to(1/3,1/3)$.  Thus all selector
claims are proved.

For the test point \eqref{eq:test}, $x^2+y^2=1$, so it lies in $\D$, and
\begin{equation}\label{eq:exactz}
 G(x)=\frac{99}{100},\qquad
 \widetilde z_2=\frac12-\frac{50}{319}
       \left(\frac{99}{100}\right)^{100},\qquad
 \widetilde z_1=\frac12-\frac9{20}\widetilde z_2.
\end{equation}
Let $C_2$ be the sum in \eqref{eq:theta} at \eqref{eq:test} and
\eqref{eq:exactz}, restricted to $-2\leq k,l\leq2$.  Positivity gives
\begin{equation}\label{eq:clower}
 \tht_{x,y}(\widetilde z;1)\geq C_2.
\end{equation}

It remains to give a rigorous upper bound for the hexagonal benchmark.  Write
$u=k+1/3$ and $v=l+1/3$.  Its exponent uses
\[
 q_H(u,v)=\frac{2}{\sqrt3}(u^2+uv+v^2)
 \geq\frac{u^2+v^2}{\sqrt3}.
\]
Set $c=\pi/\sqrt3$ and, for $N\geq0$,
\begin{align*}
 A_N&=\sum_{|n|\leq N}e^{-c(n+1/3)^2},\\
 T_N&=\frac{2e^{-c(N+2/3)^2}}
 {1-e^{-c(2N+7/3)}}.
\end{align*}

\begin{lemma}\label{lem:tail}
The complement of the square $|k|,|l|\leq N$ in the hexagonal theta sum is
at most $2(A_N+T_N)T_N$.
\end{lemma}

\begin{proof}
The one-dimensional tail satisfies
\[
 \sum_{|n|>N}e^{-c(n+1/3)^2}
 \leq2\sum_{n=N+1}^{\infty}e^{-c(n-1/3)^2}\leq T_N.
\]
Indeed, the first dominating term is $e^{-c(N+2/3)^2}$, and the ratio of
successive terms is at most $e^{-c(2N+7/3)}$.  The full one-dimensional sum
is consequently at most $A_N+T_N$.  By $q_H\geq(u^2+v^2)/\sqrt3$ and a union
bound over the two coordinate tails, the two-dimensional complement is at
most $2(A_N+T_N)T_N$.
\end{proof}

Let $H_2$ be the $5\times5$ hexagonal square.  The directed 60-decimal
interval calculation in the accompanying checker (file
\path{code/check_counterexample.py}) gives
\begin{align*}
 C_2&>0.9242317405007238239445,\\
 H_2&<0.9203713649781851694427,\\
 2(A_2+T_2)T_2&<0.000013118172843186372.
\end{align*}
Every quantity is a finite composition of rational arithmetic, square roots,
$\pi$, and exponentials, evaluated with outward rounding.  Therefore
\begin{align*}
 \tht_{x,y}(\widetilde z;1)
 -\tht_{1/2,\sqrt3/2}((1/3,1/3);1)
 &\geq C_2-H_2-2(A_2+T_2)T_2\\
 &>0.00384725734969546>0.
\end{align*}
Together with \eqref{eq:clower} and Lemma~\ref{lem:tail}, this proves the
strict reversal and completes the proof of Theorem~\ref{thm:main}.

\begin{corollary}
For the selector \eqref{eq:selector}, the conjectured inequality fails on an
open neighborhood of the parameter triple in \eqref{eq:test}.
\end{corollary}
\begin{proof}
The theta series converges locally uniformly in $(x,y,z,\alpha)$ for $y>0$.
The strict gap just proved therefore persists by continuity.
\end{proof}

\section{Verification, scope, and novelty}

The interval checker is deterministic and asserts positivity of the lower
endpoint of the certified gap.  The separate \texttt{verification.md} audits
the selector conditions, tail estimate, and possible interpretations of the
source's final boundary condition.

This counterexample refutes only the proposed sufficient conditions in
Appendix~C.  It does not affect the paper's proved variational theorem, the
Strohmer--Beaver application, or the fact that the hexagonal deep hole is a
universal cold spot.  A corrected stability statement would need a quantitative
modulus controlling how fast $\widetilde z$ may leave the circumcenter as the
lattice leaves the hexagonal boundary.

A bounded novelty search on 11 August 2026 covered the run's registry,
solution, attempt, and proof-gap indexes; exact-title and arXiv-id searches;
and web searches combining ``stability conjecture'', ``widetilde z'', and
``counterexample''.  Later work on stable cold spots cites the source theorem
and emphasizes that exact cold spots are difficult to determine
\cite{BDV}, but the search found no resolution of this Appendix selector
conjecture.  This supports, but cannot certify, novelty.

\textbf{Human-review recommendation.}  Review as a full counterexample.  The
only computational component is the explicit finite interval inequality; its
analytic tail enclosure is included above and can be independently reproduced
in under a second.

\begin{thebibliography}{9}
\bibitem{BFS}
L. B\'etermin, M. Faulhuber, and S. Steinerberger,
\emph{A variational principle for Gaussian lattice sums},
arXiv:2110.06008 (2021), Appendix~C,
\url{https://arxiv.org/abs/2110.06008}.

\bibitem{BDV}
E. Bannai, L. B\'etermin, and F. Vallentin,
\emph{Polarization of lattices: Stable cold spots and spherical designs},
Forum of Mathematics, Sigma (2026).
\end{thebibliography}

\end{document}
